% Part: lambda-calculus
% Chapter: introduction
% Section: minimization

\documentclass[../../../include/open-logic-section]{subfiles}

\begin{document}

\olfileid{lam}{int}{min}
\olsection{\usetoken{S}{lambda definable} ప్రమేయాలు కనిష్ఠీకరణ కింద సంవృతమైనవి}

\begin{lem}
$f(x,y)$ అనేది \tetoken{లాంబ్డాతో నిర్వచించదగిన}{!!{lambda definable}}దని అనుకుందాం. $g$ను
\[
g(x) \simeq \umin{y}{f(x,y)}.
\]
ద్వారా నిర్వచించండి. అప్పుడు $g$ కూడా \tetoken{లాంబ్డాతో నిర్వచించదగిన}{!!{lambda definable}}ది.
\end{lem}

\begin{proof}
స్థూలంగా భావన ఇలా ఉంటుంది. $x$ ఇచ్చినప్పుడు, $n$ వద్ద మొదలుపెట్టి
ఒక~$y$ కోసం అన్వేషించే ప్రమేయం $h_x(n)$ను నిర్వచించడానికి స్థిరబిందు
లాంబ్డా పదం~$Y$ను వాడతాం; అప్పుడు $g(x)$ కేవలం $h_x(0)$. ప్రమేయం~$h_x$ను
ఒక స్థిరబిందు సమీకరణపు పరిష్కారంగా వ్యక్తపరచవచ్చు:
\[
h_x(n) \simeq
\begin{cases}
n & \text{$f(x,n) = 0$ అయితే} \\
h_x(n+1) & \text{లేకపోతే.}
\end{cases}
\]

వివరాలు ఇవి. పరికల్పన ప్రకారం $f$ ఒక పదం~$F$చే \tetoken{లాంబ్డాతో నిర్వచించబడిన}{!!{lambda defined}}ది.
$D(M, N, \bar{0}) \red M$, $D(M, N, \bar{1}) \red N$ అయ్యే లాంబ్డా
పదం~$D$ కూడా మనకు ఉందని గుర్తుంచుకోండి. ప్రస్తుతానికి $x$ను స్థిరపరిచి,
$h_x$ను \tetoken{లాంబ్డాతో నిర్వచించే}{!!{lambda define}} ($x$పై ఆధారపడే) పదం~$H$ను కనుగొనాలంటే అది కింది సమీకరణాన్ని
నెరవేర్చాలి:
\[
H(\num{n}) \equiv D(\num{n}, H(\fn{Succ}(\num{n})), F(x, \num{n})).
\]
\sourcecorrection{OLTELAMCOMP-007}{ఉపసిద్ధాంతం $f$ లాంబ్డాతో నిర్వచించదగినదని మాత్రమే అనుకున్నప్పటికీ, ఆంగ్ల మూల నిరూపణ ఇక్కడ అది ఆదిమ పునరావృత్తమనే మరింత బలమైన, లభించని పరికల్పనను ప్రకటించింది. నిరూపణ నిజంగా వాడే పరికల్పన ప్రకారం $f$ను నిర్వచించే లాంబ్డా పదం~$F$ ఉందని నేరుగా పేర్కొన్నాం.}
\sourcecorrection{OLTELAMCOMP-008}{ఈ అన్వేషణ నిర్మాణంలో ఆంగ్ల మూలం నిర్వచించని ఉత్తరగామి పదం~$S$ను నాలుగు సార్లు వాడింది. మునుపటి ఉపసిద్ధాంతం నిర్వచించిన $\fn{Succ}$ పదంతో ప్రతిసారీ దాన్ని పునరుద్ధరించాం.}
స్థిరబిందు పదం~$Y$ను వాడి దీన్ని చేయవచ్చు. మొదట $U$ను కింది పదంగా
తీసుకోండి:
\[
\lambd[h][\lambd[z][D(z,(h(\fn{Succ}z)),F(x,z))]],
\]
తరువాత $H$ను $YU$ అనే పదంగా తీసుకోండి. $H$లోని ఏకైక స్వేచ్ఛా చరం~$x$
అని గమనించండి. $H$ పైనున్న సమీకరణాన్ని నెరవేరుస్తుందని చూపుదాం.

$Y$ నిర్వచనం ప్రకారం,
\[
H = YU \equiv U(YU) = U(H).
\]
ప్రత్యేకంగా, ప్రతి సహజ సంఖ్య~$n$కు,
\begin{align*}
H(\num{n}) & \equiv U(H, \num{n}) \\
& \red D(\num{n}, H(\fn{Succ}(\num{n})), F(x, \num{n})),
\end{align*}
అవుతుంది; మనకు కావలసింది ఇదే. చివరి పంక్తిలో $x$ స్థానంలో సంఖ్యాంకం
$\num{m}$ను ప్రతిస్థాపిస్తే, $F(\num{m}, \num{n})$ అనేది $\num{0}$కు
తగ్గినప్పుడు ఆ వ్యక్తీకరణ $\num{n}$కు తగ్గుతుంది; $F(\num{m},
\num{n})$ మరే సంఖ్యాంకానికి తగ్గితే $H(\fn{Succ}(\num{n}))$కు తగ్గుతుంది.

నిరూపణను ముగించడానికి, $G$ను $\lambd[x][H(\num 0)]$గా తీసుకోండి.
అప్పుడు $G$ ప్రమేయం~$g$ను \tetoken{లాంబ్డాతో నిర్వచిస్తుంది}{!!{lambda define}s}; మరో మాటలో, ప్రతి~$m$కు
$g(m)$ నిర్వచితమైతే $G(\num m)$ అనేది~$\overline {g(m)}$కు తగ్గుతుంది;
లేకపోతే దానికి నియత రూపం ఉండదు.
\end{proof}

\end{document}
